\subsection{A quantitative moment for the prescribed seed}
\label{global-gram-quantitative-seed}

The qualitative seed theorem contains an uncontrolled reciprocal
compact mass.  We remove it by estimating the actual normalized
seed directly.  The additional tensor steps below are proved here;
they are not claimed to lie in the earlier numerical family
$q_{\rm fam}$.  The original arithmetic representations, periods,
finite columns and Haar measures are retained.

First work at one real place.  Put $G=U(p,q)$, where
$p=k+r$, $q=r+1$, $k\ge1$, $n=p+q\ge1801$, and $m=2r+2$.
Use a Euclidean orthonormal Witt basis with
$h(e_a,f_b)=\delta_{ab}$ and a central vector
$z=\sqrt T\,e_-$ of Hermitian norm $-T$, where $T\ge1$.
Write $S=N\rtimes U(k)$ for the flag Bessel subgroup.  Its
character is the literal standard character
\[
 \psi(n)=\exp\bigl(2i\operatorname{Im}\ell(n)\bigr),\qquad
 \ell(n)=\sum_{a<r}h(f_a,n e_{a+1}),\quad e_r:=z.
\]
Indeed the source defines $c_{\rm Std}=-2i$ and
$\psi_c=\exp(i\operatorname{Re}(c\ell))$.
The empty sum is zero when $r=0$.
Let $\sigma$ be the actual discrete series of integral
Harish--Chandra parameter $\lambda$, let $w$ be any unit vector
in its lowest compact type, and retain the actual compact target
of parameter $u=(u_1,\ldots,u_k)$.  Assume that every noncompact
root gap is at least $16n^{19}$.  Set
\begin{equation}\label{global-gram-quantitative-B}
 B=10^6n^7+1+\max_a|\lambda_a|+\max_i|u_i|,
 \qquad H=100n^2B,
 \qquad K=T H^{5n}.
\end{equation}
All group and subgroup measures are the given measures; their
quotient $dq$ is defined by the Weil formula.
\cite{global-gram-quantitative-source}

\paragraph{The additional dual tensor family.}
Starting at $\lambda$, construct the nonzero factors of every
path of at most $M=10n^2$ successive dual-standard tensor
products.  A path replaces $\mu$ by $\mu-e_s$ at each step.
Every noncompact gap stays at least
$16n^{19}-10n^2\ge8n^{19}$.  An integral ordered parameter can
leave its open chamber in a single such step only by an adjacent
equality.  That root must be compact.  The corresponding induced
factor is zero: induction in stages through its compact rank-one
Levi gives $\mathcal O_{\mathbb P^1}(-1)$, whose cohomology is
zero in both degrees.  This proves the zero-factor convention,
including when an earlier compact gap equals one.

All other factors are regular discrete series in the same chamber.
The tensor identity and good-range vanishing give the finite
filtration, and its distinct one-step central characters split it.
Casselman--Wallach globalization and the invariant Hermitian form
give smooth intertwiners $i_s,p_s$ with
\[
 p_si_t=\delta_{st},\qquad \sum_s i_sp_s=1,\qquad
 [i_sv,i_sv]=\epsilon_s\|v\|^2,\quad \epsilon_s\in\{1,-1\}.
\]
For a zero factor set $i_s=p_s=0$; the sum runs over the
nonzero factors.  Iteration keeps path labels, so repeated final
parameters never require a self-extension splitting assertion.
This is the proof of the established stable tensor lemma with
the dual weights $-e_s$ in place of $+e_s$, using its cited
Harish--Chandra, cohomological-induction and globalization
theorems.\cite{current-construction-record}

The normalized Bessel-family theorem in ``Constructing the local
representation and period families'' applies to every one of
these auxiliary representations.\cite{current-construction-record}
Its convergence proof
uses the noncompact gap, not compact cloud separation or target
interlacing.  The gap $8n^{19}$ satisfies its displayed numerical
bounds.  Amenability and Bessel multiplicity one then give a
possibly zero map $\mathcal B_\mu$, normalized by the unchanged
subgroup period.  The same quotient-mass theorem gives
\begin{equation}\label{global-gram-quantitative-mass}
 F_\mu(v;g)=\sqrt{d_\mu}\,\mathcal B_\mu\sigma_\mu(g)v,
 \qquad \int\|F_\mu(v;g)\|^2dq(g)\le\|v\|^2.
\end{equation}
If the map is nonzero this is an equality.  No nonvanishing is
assumed for the auxiliary maps.

\begin{lemma}[Coarse projection bound with the larger reserve]
\label{global-gram-quantitative-projection}
Start with the whole lowest compact type at $\lambda$.  Permit
at most $10n^2$ successive tensor projections and at most
$6n^3$ Lie-algebra letters, in any order for which the path exists.
For every resulting vector $v$, whenever another tensor step is
permitted, and every positive orthonormal dual coordinate $e_c$,
\[
       \sum_s\|p_s(e_c\otimes v)\|^2\le H^{2n}\|v\|^2.
\]
On all the compact modules needed in this argument,
$\sum_{a,b}\|E_{ab}v\|^2\le H^2\|v\|^2$.
\end{lemma}
\begin{proof}
The lowest compact Casimir has square-root bound $2\sqrt nB$:
its highest weight is $\Lambda=\lambda+\rho-2\rho_K$, so its
Casimir is $\|\lambda+\rho-\rho_K\|^2-\|\rho_K\|^2$.
Every auxiliary parameter has absolute value at most $B$.
Tensoring with a standard or dual coordinate adds at most
$\sqrt n$ to the square root of a compact energy bound.
A Lie letter adds at most $2n$, by the Hilbert--Schmidt adjoint
Casimir calculation in ``Simultaneous moments of flag-length
ratios''.\cite{current-construction-record}
A compact intertwiner introduces no new compact Casimir
eigenvalues; this assertion does not require a Hilbert isometry.
Applying an inclusion and then taking a coordinate costs at most
two extra $\sqrt n$ terms.  Thus all components used below have
compact energy radius at most
\[
 2\sqrt nB+(10n^2+2)\sqrt n+2n(6n^3+2)\le4nB.
\]
The full Casimir is $c_\mu=\|\mu\|^2-\|\rho\|^2\ge-n^3$.
Consequently the exact energy identity gives
\begin{equation}\label{global-gram-quantitative-energy}
 \sum_{a,b}\|E_{ab}v\|^2
 =2\langle v,\Omega_Kv\rangle-c_\mu\|v\|^2
 \le33n^2B^2\|v\|^2\le(6nB)^2\|v\|^2.
\end{equation}

Put $P_s=i_sp_s$ and $(\mathcal Tx)_a=-\sum_bE_{ab}x_b$.
For $x=e_c\otimes v$, let $V_x$ be the finite span of $P_sx$.
The tensor Casimir identity gives
$\mathcal TP_s=\zeta_sP_s$, where
$\zeta_s=-\mu_s+(1-n)/2$.
The components of $V_x$ satisfy the bound just proved.
Cauchy--Schwarz in each row gives
$\|\mathcal T|_{V_x}\|\le6n^2B$.
Furthermore $|\zeta_s|\le2B$ and
$|\zeta_s-\zeta_t|=|\mu_s-\mu_t|\ge1$ for $s\ne t$.
Polynomial interpolation, including the harmless zero-factor
labels, therefore gives
\[
 P_s|_{V_x}=\prod_{t\ne s}
       \frac{\mathcal T-\zeta_t}{\zeta_s-\zeta_t},
 \qquad \|P_s|_{V_x}\|\le H^{n-1}.
\]
This is a statement on a finite invariant space, so no bound for
an unbounded Lie operator on the ambient Hilbert space is used.
Finally
\[
 \sum_s\|p_sx\|^2
 =\sum_s|[P_sx,P_sx]|
 \le nH^{2n-2}\|x\|^2\le H^{2n}\|v\|^2.
\]
Equation~\eqref{global-gram-quantitative-energy} also proves the
asserted derivative bound, since $6nB\le H$.
\end{proof}

\paragraph{Lowering constants without an interlacing reserve.}
For every nonzero base period map and every regular shifted
factor, the exact cofactor evaluation and period-adjoint argument
in ``The degree-normalized Bessel lowering scalar'' give
\cite{current-construction-record}
\begin{equation}\label{global-gram-quantitative-lowering}
 |\kappa_s|^2\frac{d_\mu}{d_{\mu-e_s}}
 =T\frac{\prod_{i=1}^k|\mu_s-1/2-u_i|}
          {\prod_{t\ne s}|\mu_s-1-\mu_t|}
 \le T H^n.
\end{equation}
The denominator factors are nonzero integers and have absolute
value at least one.  The numerator factors are at most $3B$.
These facts replace the earlier finer interlacing estimate.
The actual root character has the opposite sign to the adjacent
imaginary-coordinate convention in the earlier ratio display;
the Capelli root factors still have moduli $1$ and $\sqrt T$,
so this absolute scalar identity is unchanged.  If a target
period map is zero its scalar is zero.  If a compact-wall factor
is zero, omit it before applying the displayed formal-degree
ratio; division by its zero denominator is never made.
If the base map is zero all its moments vanish directly.

\begin{lemma}[High flag moments]
\label{global-gram-quantitative-flag-moments}
Let $L_i(g)$ be the Euclidean Gram--Schmidt lengths of
$g^{-1}e_0,\ldots,g^{-1}e_{r-1},g^{-1}z$, and put
$R_i=L_{i+1}/L_i$.  For the lowest-type unit vector $w$,
\[
 \int L_0^{2a}\|F_\lambda(w)\|^2dq\le K^a
       \quad(0\le a\le10n^2),
 \qquad
 \int\prod_{i<r}R_i^{2\alpha_i}\|F_\lambda(w)\|^2dq
       \le H^{2|\alpha|}\quad(|\alpha|\le6n^3).
\]
\end{lemma}
\begin{proof}
The exact dual-tensor transfer identity in the proof of
``Moments of the first flag length'',\cite{current-construction-record} with
\eqref{global-gram-quantitative-lowering}, gives
\[
 L_0^2\|F_\mu(v)\|^2
 \le nT H^n\sum_{d,s}
           \|F_{\mu-e_s}(p_s(e_d\otimes v))\|^2.
\]
At each induction step the sum of squared vector norms is at most
$nH^{2n}\|v\|^2$.  Thus the cost is
$n^2TH^{3n}\le TH^{5n}=K$.
Use \eqref{global-gram-quantitative-mass} for the zero-order
case and induct, decreasing the remaining tensor depth at every
step.  All vectors stay in the larger reserve just proved.
Tonelli justifies the induction without assuming a weighted
integral finite in advance.

For completeness, the rank-two adjacent-root vector in the proof
of ``Simultaneous moments of flag-length ratios''
\cite{current-construction-record} differentiates the present
character to $+2iL_{i+1}^2$, instead of $-2iL_{i+1}^2$.
Taking squared norms gives the same pointwise inequality
\[
 R_i^2\|F_v\|^2\le\sum_{a,b}\|F_{E_{ab}v}\|^2.
\]
Induct on $|\alpha|$, using
\eqref{global-gram-quantitative-energy} at every derivative
step and the mass bound at the last step.  This proves the
second assertion through the stated $6n^3$ reserve.
\end{proof}

\begin{lemma}[The quotient norm in flag coordinates]
\label{global-gram-quantitative-quotient-geometry}
For $\nu(g)=\inf_{s\in S}\|sg\|_{\rm op}$ one has
\begin{equation}\label{global-gram-quantitative-nu}
 1\le\nu(g)\le
 \max\left(\{1,\sqrt2L_r(g)/\sqrt T\}\cup
                  \{L_i(g),L_i(g)^{-1}:i<r\}\right).
\end{equation}
\end{lemma}
\begin{proof}
Gram--Schmidt subtraction on the ordered flag is upper
unitriangular.  It preserves its Hermitian Gram matrix and
extends by Witt's lemma to an element of $S$.
The residual isotropic vectors are Euclidean orthogonal and
Hermitian orthogonal.  Their positive and negative components
are consequently two orthogonal systems with lengths
$L_i/\sqrt2$.  A maximal compact element sends their normalized
vectors to the standard Witt vectors.

The last residual is Euclidean and Hermitian orthogonal to this
isotropic system and has Hermitian norm $-T$.  Its two component
lengths are
\[
 \alpha=\sqrt{(L_r^2-T)/2},\qquad
 \beta=\sqrt{(L_r^2+T)/2}.
\]
In particular $L_r\ge\sqrt T$.  The remaining compact group
moves it to one positive and one negative coordinate; $k\ge1$
ensures the positive coordinate is available.  It is the image
of $z$ by a boost with
$\sinh t=\alpha/\sqrt T$, $\cosh t=\beta/\sqrt T$.
Its operator norm is
$e^t=(\alpha+\beta)/\sqrt T\le\sqrt2L_r/\sqrt T$.

Together with the boosts
$e_i\mapsto L_i e_i$, $f_i\mapsto L_i^{-1}f_i$, this constructs
a representative in $g^{-1}S$ of operator norm at most the
displayed maximum.  Indeed it and the Gram--Schmidt-adjusted
$g^{-1}$ agree on the full flag, and their quotient fixes that
flag pointwise, hence belongs to $S$.
An indefinite unitary matrix and its inverse have the same
operator norm.  Inverting this representative proves the upper
bound in \eqref{global-gram-quantitative-nu}; the lower bound
follows from the determinant having modulus one.
The argument also covers $r=0$ using only the last boost.
\end{proof}

\begin{theorem}[An explicit normalized quotient moment]
\label{global-gram-quantitative-moment}
With the exact period normalization above,
\begin{equation}\label{global-gram-quantitative-moment-bound}
 \int\nu(g)^{10n^2}\|F_\lambda(w;g)\|^2dq(g)
       \le m\bigl(2T H^{6n}\bigr)^{5n^2}.
\end{equation}
This bound applies to the actual unit seed fixed by its original
matrix-coefficient projector.
\end{theorem}
\begin{proof}
For $a=5n^2$ and $0\le i\le r$, write
$L_i=L_0\prod_{j<i}R_j$ and apply Cauchy--Schwarz to the
measure $\|F_\lambda(w)\|^2dq$.  The first factor requires the
first-length order $2a=10n^2$, and the second the ratio order
$2ai\le10n^2r\le5n^3$.  Lemma~\ref{global-gram-quantitative-flag-moments}
therefore gives
\[
 \int L_i^{2a}\|F_\lambda(w)\|^2dq\le K^aH^{2ai}.
\]
Also $L_i^{-1}=L_r^{-1}\prod_{j=i}^{r-1}R_j$, so
\[
 \int L_i^{-2a}\|F_\lambda(w)\|^2dq
       \le T^{-a}H^{2a(r-i)}\quad(i<r).
\]
Raise \eqref{global-gram-quantitative-nu} to $2a$ and bound the
maximum by the sum of its $2r+2$ entries.  Since $T\ge1$,
$K=TH^{5n}$ and $5n+2r\le6n$, every resulting integral is at
most $(2TH^{6n})^a$.  This proves the first assertion.

For the last assertion let
$p(g)=d_\lambda\overline{\langle w,\sigma(g)w\rangle}$.
The original gap makes $p$ integrable, and $R(p)$ is the
orthogonal projection from the quotient-mass proof.
Its kernel on compactly supported bounded sections is exactly
$F_\lambda(w;g)F_\lambda(w;h)^*$, with the original measures.
If $\mathcal B_\lambda=0$ this kernel is zero, so density gives
$R(p)=0$.  A prescribed unit seed $W$ with $R(p)W=W$ excludes
that case.  The mass isometry and the same kernel show that
$R(p)$ has range $\mathbb C F_\lambda(w)$.
Consequently $W$ differs from $F_\lambda(w)$ by a unit scalar.
The displayed estimate is therefore for the prescribed seed;
it makes no additional Bessel nonvanishing assumption.
\end{proof}

\paragraph{The original quotient measure in the $L^1$ bound.}
Define $c_G>0$ by the exact Cartan integration formula for $dg$,
using probability maximal-compact measures and radial density
\[
 J(t)=\prod_{i<j}\sinh(t_i-t_j)^2\sinh(t_i+t_j)^2
       \prod_i\sinh(t_i)^{2(p-q)}\sinh(2t_i),
 \quad t_1\ge\cdots\ge t_q\ge0.
\]
Any fixed root-coordinate factor is included in $c_G$.
Put $s_0=\mu_S\{s:\|s\|_{\rm op}\le2\}>0$.
This set is compact and contains a neighborhood of the identity.
Let $d_0=2pq+1$.  Since
$J(t)\le\exp(2\sum_i(n-2i+1)t_i)$, integration over
$t_1\le\log R$ gives
\[
 \mu_G\{g:\|g\|_{\rm op}\le R\}\le c_G R^{d_0}
 \quad(R\ge1);
\]
here $(\log R)^q/q!\le R$.  A quotient point with $\nu\le R$
has a representative of norm at most $2R$.  Its translates by
the subgroup ball of norm two lie in the group ball of norm
$4R$.  The Weil formula thus gives
\[
 \mu_Q\{\nu\le R\}\le(c_G/s_0)4^{d_0}R^{d_0},\qquad
 \int_Q\nu^{-2n^2}dq\le2(c_G/s_0)4^{d_0}.
\]
The last inequality is layer cake and
$2n^2/(2n^2-d_0)\le2$.
Cauchy--Schwarz with
\eqref{global-gram-quantitative-moment-bound} proves the actual
weighted bound
\begin{equation}\label{global-gram-quantitative-L1}
 \int_Q\nu^{4n^2}\|W\|dq
 \le {\cal L}(T,B;c_G,s_0):=
 \sqrt{\frac{2m c_G}{s_0}}\,2^{2pq+1}
                    (2T H^{6n})^{5n^2/2}.
\end{equation}
The ratio $c_G/s_0$ is retained.  Rescaling $ds$ changes the
normalized seed and this $L^1$ norm, so it cannot be silently
replaced by one.

\paragraph{Exact transport from the original arithmetic frame.}
At every real place use the source coordinates $P_v^{-1}$.
Let $t_{{\rm fr},v}\in G_v$ carry the standard Witt frame, the
vector $\sqrt{T_u}e_-$, and its positive orthonormal complement
to the actual $e_a,f_a,z_u,w_{u,v}$.  Such an isometry exists
because their Hermitian Gram matrices agree.  Put
\begin{equation}\label{global-gram-quantitative-character-scale}
 c=\frac{g}{2\pi},\qquad
 \eta_a=(-ic)^{r-a},\qquad
 D_\eta e_a=\eta_a e_a,\quad
 D_\eta f_a=\overline{\eta_a}^{-1}f_a,\quad
 D_\eta z=z,\qquad t_v=t_{{\rm fr},v}D_\eta.
\end{equation}
Here $g$ is the actual product of the two selected rational
primes; it is the same scalar at every embedding.
Every adjacent root coordinate, including the last one involving
$z$, is multiplied under conjugation by $D_\eta$ by $-ic$.
The literal arithmetic character is
$\exp((4\pi i/g)\operatorname{Re}\sum_a B_a)$ at each place.
Thus
\[
 \psi_{u,v}(t_vnt_v^{-1})
 =\exp\bigl((4\pi i/g)\operatorname{Re}(-ic\ell(n))\bigr)
 =\exp(2i\operatorname{Im}\ell(n)).
\]
Multiplying these equalities proves the exact global character
identity; no degree or embedding factor is removed.
The diagonal $D_\eta$ is the identity on the positive compact
complement, so the compact target remains the prescribed one.

Set $\widetilde W_u(g)=W_u(tg)$, push $ds_u$ by
$s\mapsto t^{-1}st$, and push the given representative measure
$\nu_Q(u)$ by $x\mapsto t^{-1}x$.  Left invariance keeps $dg$
unchanged, and the Weil formula gives exactly the quotient
measure used above.  Right convolution by the actual $p$ is
unchanged.  Haar uniqueness allows the given product-group
subgroup measure to be factored into local Haar measures whose
product is exactly that measure; no total scalar is changed.

The source's actual tensor identity for $W_u$ then gives local
continuous sections.  Its global equivariance implies their
local equivariance by fixing nonzero factors at the other
places.  Its unit norm and Tonelli give finite positive local
norms $a_v$ with $\prod_v a_v=1$.  The product projector fixes
their tensor, hence each normalized local section is fixed by
its local orthogonal projector.  The last part of
Theorem~\ref{global-gram-quantitative-moment} applies to each.
Since the norm is the product of local norms, the resulting
global estimate is exactly
\begin{equation}\label{global-gram-quantitative-global-L1}
 \int\left(\prod_v\nu_v(g_v)\right)^{4n^2}
           \|\widetilde W_u(g)\|\,dq(g)
 \le\prod_{v\mid\infty}{\cal L}(T_u,B_v;c_{G,v},s_{0,v}).
\end{equation}
The infimum for the product quotient norm equals the product
of the local infima, since the subgroup is a direct product.
All $f=[F:\mathbb Q]$ factors remain.

\paragraph{The exact weighted far-return consequence.}
The arithmetic minimum proof can now be applied at the original
point $(tg,a_{u,\rho}k)$ using the vectors $t_v^{-1}y_v$.
This is an equality of its actual metric realization
$(t_vg_v)^{-1}y_v=g_v^{-1}t_v^{-1}y_v$, not a change of the
arithmetic lattice.  The transformed norms of its flag vectors
are $c^{-(r-a)}$, $c^{r-a}$ and $\sqrt{T_u}$.
The four arithmetic pairing cases and every higher-rank
covolume normalization therefore give
\[
 \theta_u^{\rm st}=
 \min\bigl(\{1,T_u^{-1/2}\}\cup
 \{c^{r-a},D_F^{-(r-a)/f}c^{-(r-a)}:a<r\}\bigr),\qquad
 \tau_u^{\rm st}=\min\left\{\tfrac12,
       \frac{\theta_u^{\rm st}}{\sqrt nD_E^{1/(4f)}}\right\},
\]
and
\[
 \tau(tg,a_{u,\rho}k)\ge
           \tau_u^{\rm st}\left(\prod_v\nu_v(g_v)\right)^{-1/f}.
\]
The fractional different is unchanged: its absolute $E$-norm
is still $D_F^2$.  In particular no $f_a$ is declared integral.

Under precisely the cited \texttt{Deep.FarBound} hypotheses
in Section~\ref{global-gram-arithmetic-minimum}, define
\[
 C_u^{\rm st}=2f3^f
        \left(\frac{3\sqrt n}{(\tau_u^{\rm st})^2}\right)^{2n^2f}.
\]
The original normalized far-return sum $\mathcal T_{h,\varepsilon}$
then satisfies, for every original column and every $k\in K_2$,
\begin{equation}\label{global-gram-quantitative-far}
 \int \mathcal T_{h,\varepsilon}(tg,a_{u,\rho}k)
         \|\widetilde W_u(g)\|\,d(t^{-1}_*\nu_Q(u))(g)
 \le C_u^{\rm st}e^{-h\varepsilon^2/16}
       \prod_v {\cal L}(T_u,B_v;c_{G,v},s_{0,v}).
\end{equation}
The original normalizer is
$(\prod_vd_v)e(1)$, with $e(1)=\operatorname{vol}(K_2)^{-1}$.
The far-return factor itself need not descend to the $S$ quotient;
the displayed representative measure is consequently retained.
Its majorant does descend, which is exactly what
\eqref{global-gram-quantitative-global-L1} uses.

This is a quantitative bound for the actual seed and its actual
far-return integral.  It replaces the unknown compact-mass
reciprocal and unknown frame heights by the displayed arithmetic,
weight and Haar quantities.  Bounds for those quantities across
the moving family, the near returns, and the exact retained
arithmetic-chain trace remain required for W08.  The cited
\texttt{Deep.FarBound} is not reproved here, and no final numerical
trace budget follows merely from this estimate.
